% !TeX root = ../collection-solutions.tex

\titledquestion{Falling behind: inequality and household debt}

\citet{drechsel2026falling} ask whether rising income inequality fueled the American household debt boom between 1980 and 2007.

\begin{parts}
    \part[4] What is the key takeaway of the paper? Explain how the authors arrive at this conclusion: what is the central theoretical condition, and how is the quantitative experiment set up?

    \begin{solutionorbox}[12cm]
        \emph{Takeaway} (2 points): Rising income inequality raises aggregate housing demand, house prices and mortgage debt \emph{if and only if} social comparisons are asymmetric — households look \emph{up} to the rich rather than to the average. With the standard ``keeping up with the (mean) Joneses'' specification, changes in the income distribution have exactly \emph{zero} aggregate effect. Quantitatively, the observed rise of the top-10\% income share (from $33\%$ to $41\%$) combined with upward-looking comparisons accounts for the bulk ($\approx 75\%$) of the rise in the housing expenditure share, and a meaningful share ($\approx 10$--$15\%$) of the mortgage and house price booms.

        \emph{How they get there} (2 points):
        \begin{itemize}
            \item Theory: a tractable dynamic model where housing yields status relative to a reference house determined by a comparison network. A redistribution toward a group raises aggregate housing and debt if and only if that group is more ``popular'' (more looked-at) than the group it gains from; comparison to the mean makes everyone equally popular, so aggregates are invariant to the distribution.
            \item Quantification: calibrate to the US in 1980 (mortgage-to-income, housing expenditure share); discipline the strength of the comparison motive with an external micro estimate (Bellet's housing-satisfaction evidence); then feed in the observed shift of the income distribution toward the top 10\% as the only shock.
            \item The model also reproduces \emph{who} took on debt — debt-to-income rises for the bottom 50\% and middle 40\%, not for the rich — and why wealth inequality rose by less than income inequality.
        \end{itemize}
        Grading: 1 point for the if-and-only-if role of upward-looking comparisons (incl.\ zero effect under mean comparisons), 1 point for the quantitative bottom line; 1 point for the calibration/external-discipline setup, 1 point for feeding in the observed income shift as the only shock.
    \end{solutionorbox}

    \part[4] Explain intuitively how an increase in \emph{top} incomes can raise the mortgage debt of \emph{non-rich} households. Make clear (i) why the response is driven entirely by the non-rich, and (ii) why the durability of housing is essential for the \emph{debt} (rather than just the housing) response.

    \begin{solutionorbox}[12cm]
        \begin{itemize}
            \item (1 point) When top incomes rise, the rich buy bigger houses. For everyone who compares their house to the houses of the rich, the reference level rises and their housing \emph{status} falls. Non-rich households respond by shifting expenditure from (status-neutral) consumption toward (status-relevant) housing — they keep up with the rich Joneses.
            \item (1 point) The rich themselves do not change their expenditure shares: they are at the top of the comparison network, so their reference point does not move with others' choices, and with homothetic (Cobb--Douglas) demand their housing simply scales with income. The entire adjustment therefore comes from the non-rich — matching the data, where debt-to-income rose for the bottom 90\% but not for the top 10\%.
            \item (2 points) Housing is \emph{durable}: the cheapest way to finance a permanently larger house is to buy it upfront with a mortgage and spread the cost over future income. A bigger target house therefore translates directly into more mortgage debt. If housing were non-durable, housing \emph{spending} would still rise but debt would not respond. Consistent with this, rising top incomes predict higher non-rich \emph{mortgage} debt in the data, but not higher non-mortgage (consumption) debt.
        \end{itemize}
    \end{solutionorbox}
\end{parts}
